\bookchapter{Chapitre 68 : Le chemin de l'avancement — permissions, notifications et architecture moderne}

\chapterintro{Savoir fabriquer une App n'est que le commencement. Ce chapitre discute des marches à franchir entre « savoir faire » et « bien faire » : demander les permissions avec élégance, envoyer des notifications sans devenir insupportable, élever l'architecture d'un cran encore, et — en 2026 — écrire du code iOS en tandem avec l'IA. Chacun de ces sujets se présentera à toi tout de suite dans un vrai projet.}

\booksection{Les permissions : frapper d'abord, entrer ensuite}

Caméra, localisation, contacts, notifications… pour ces « zones sensibles », iOS exige que l'App demande à l'utilisateur par une fenêtre au moment de l'usage. La règle tient en trois pas : « déclarer + demander + dégrader avec élégance » :

\begin{lstlisting}[language=Swift]
// 1. déclare dans Info.plist le texte d'usage (sans déclaration, crash immédiat !)
// Privacy - Location When In Use Usage Description: "Servira à consulter la météo de ta ville"

// 2. la demander au moment de l'usage
import CoreLocation

@Observable
class LocationModel {
    private let manager = CLLocationManager()

    func request() {
        switch manager.authorizationStatus {
        case .notDetermined: manager.requestWhenInUseAuthorization()
        case .denied, .restricted: /* guider vers les réglages pour l'activer */
        default: /* accordé, on y va */
        }
    }
}

// 3. dégrader élégamment après un refus : afficher un message et aller aux réglages
// UIApplication.shared.open(URL(string: UIApplication.openSettingsURLString)!)
\end{lstlisting}

\begin{tcbwarn}{Attention : les permissions iOS ne passent qu'« une fois »}
La fenêtre de permission iOS ne s'affiche qu'une seule fois — après un refus, elle ne reviendra jamais. Donc : choisis le moment de la demande « à l'instant où l'utilisateur s'apprête à s'en servir » (et non une salve au lancement), et écris un texte qui explique l'usage. Le guidage d'interface après refus (« active la localisation dans Réglages ») est d'équipement ; ne laisse pas l'utilisateur enfermé dans une impasse.

\end{tcbwarn}

\booksection{Les notifications : dire des choses utiles au bon moment}

Les notifications sont le canal par lequel l'App rejoint l'utilisateur, et la fonction qu'il « bloque » le plus volontiers. Avec le framework UserNotifications d'iOS, trois disciplines pour en envoyer : ne diffuser que l'utile, laisser l'utilisateur couper, ne pas harceler au milieu de la nuit. L'usage de base :

\begin{lstlisting}[language=Swift]
import UserNotifications

// 1. demander la permission de notification (une seule fenêtre)
UNUserNotificationCenter.current()
    .requestAuthorization(options: [.alert, .sound, .badge]) { granted, _ in
        print(granted ? "Accordé" : "Refusé")
    }

// 2. programmer un "rappel de dépenses chaque jour à 20 h"
func scheduleDailyReminder() {
    let content = UNMutableNotificationContent()
    content.title = "Rappel de dépenses"
    content.body = "Tu n'as encore rien noté aujourd'hui"
    content.sound = .default

    var comps = DateComponents()
    comps.hour = 20
    let trigger = UNCalendarNotificationTrigger(dateMatching: comps, repeats: true)

    let request = UNNotificationRequest(
        identifier: "daily-reminder",
        content: content,
        trigger: trigger
    )
    UNUserNotificationCenter.current().add(request)
}
\end{lstlisting}

Attention : la permission de notification d'iOS, elle aussi, ne se demande « qu'une fois », et l'utilisateur peut la couper à tout instant dans les réglages du système — la valeur perçue du contenu des notifications est donc capitale : un seul abus peut te priver de l'utilisateur pour toujours.

\booksection{Architecture avancée : MVVM en profondeur et testabilité}

Tu as déjà utilisé MVVM au chapitre 65 ; voici trois points avancés pour compléter :

\begin{itemize}
  \item L'abstraction par protocole : fais dépendre le ViewModel d'un protocole (protocol) plutôt que d'une classe concrète — aux tests, tu peux injecter une « fausse source de données » ; c'est la pierre d'angle de la testabilité ;
  \item Responsabilité unique : un ViewModel ne s'occupe que d'une seule chose (celui des comptes ne gère pas la météo) ; quand les fonctions s'entassent, on scinde ;
  \item Ne surconçois pas : pour un petit projet en solo, « classe @Observable + abstraction par protocole » suffit amplement. Les architectures lourdes de type VIPER ou Clean Architecture attendront que l'équipe et la complexité sonnent à ta porte — les outils servent un objectif, ils ne servent pas à se pavaner.
\end{itemize}

\begin{tcbinfo}{Astuce : les « dividendes de l'écosystème » de SwiftUI}
L'écosystème iOS cache encore deux façons avancées « réservées à la maison » : les widgets (WidgetKit) laissent l'utilisateur lire l'information sans ouvrir l'App (les « dépenses du mois » du carnet posées sur l'écran d'accueil), et les raccourcis (App Intents) laissent l'utilisateur invoquer tes fonctions par Siri ou l'automatisation. Les deux s'appuient sur des frameworks officiels, peu coûteux à apprivoiser, mais au gain d'expérience utilisateur énorme — sur la route de l'avancement, ce sont des déverrouillages à traiter en priorité.

\end{tcbinfo}

\booksection{Les nouveautés de 2026 : le développement iOS assisté par l'IA}

En 2026, l'IA est déjà profondément enchâssée dans le flux de développement Apple. Si tu veux gagner en vitesse, cette combinaison est bien rodée :

\begin{center}
\begin{tabularx}{\linewidth}{|X|X|X|}
\hline
\textbf{Étape} & \textbf{Comment l'IA aide} & \textbf{Ta responsabilité} \tabularnewline
\hline
Écrire le code & Complétion de code intégrée à Xcode + génération par des plugins IA tiers & Décrire clairement le besoin, passer au crible le produit \tabularnewline
Chercher la documentation & Donner la doc officielle et les messages d'erreur à l'IA pour les décoder & Revenir à la doc officielle pour vérifier les conclusions clés \tabularnewline
Écrire les tests & Faire générer par l'IA des cas XCTest en trois classes « normal / limites / exception » & Confirmer que les cas collent au sens métier \tabularnewline
Corriger les bugs & Coller le journal de plantage + le code concerné, l'IA localise et répare & Comprendre les modifications, lancer les tests pour valider \tabularnewline
\hline
\end{tabularx}
\end{center}

\begin{center}\small\color{gray}Les quatre scénarios du développement iOS assisté par l'IA\end{center}

La méthode concrète — comment donner des instructions à l'IA, comment passer la recette, comment éviter les pièges — occupe onze chapitres complets dans la septième partie de ce livre, « Pratique de la programmation avec l'IA » (chapitres 37-47) ; je te conseille chaudement de lire les deux en parallèle : la syntaxe de Swift pour cette partie, l'art de collaborer avec l'IA pour la septième.

\begin{tcbpractice}{À toi de jouer : ajoute le « trio » au carnet de comptes}
\booksection{L'essentiel de la concurrence Swift : @MainActor et Task}

Le régime de concurrence de Swift 6 comporte deux règles que tout débutant doit comprendre :

\begin{itemize}
  \item L'interface ne se met à jour que sur le thread principal : les vues SwiftUI sont @MainActor par défaut ; les rappels réseau reviennent sur le thread principal avant de toucher à l'état (await MainActor.run \{ ... \}, ou bien marque le ViewModel comme @MainActor) ;
  \item Les tâches asynchrones se pilotent avec Task : le modificateur .task \{ \} démarre tout seul quand la vue apparaît et s'annule tout seul à sa disparition — c'est la posture standard pour lancer un travail asynchrone en SwiftUI ; évite d'écrire des Task à nu en vrac dans le body.
\end{itemize}

\begin{lstlisting}[language=Swift]
struct WeatherScreen: View {
    @State private var vm = WeatherViewModel()

    var body: some View {
        Group {
            switch vm.state {
            case .loading: ProgressView()
            case .success(let d): WeatherContent(data: d)
            case .error(let m): Text(m)
            }
        }
        .task {                       // démarre à l'apparition, s'annule à la disparition
            await vm.search(city: "Paris")
        }
    }
}
\end{lstlisting}

Derrière ces deux règles se tient la conception de Swift 6 « sûr contre les courses de données » — le compilateur te signale l'erreur dès que tu écris de travers, au lieu d'attendre que l'App plante pour le découvrir. Les erreurs du compilateur ne te cherchent pas des noises : elles interceptent les balles à ta place.

(1) Ajoute au carnet le « rappel de comptes quotidien à 20 h » (permission de notification + notification locale) ; (2) transforme la source de données du ViewModel en abstraction par protocole et écris un test unitaire qui injecte de fausses données ; (3) ajoute un widget d'écran d'accueil « dépenses du mois ». Le trio terminé, ton projet aura les allures d'une « troupe régulière ».

\end{tcbpractice}

\begin{tcbsummary}{Résumé du chapitre}
\begin{itemize}
  \item Trois pas de la permission : déclaration dans Info.plist (l'omettre = plantage assuré), demande au moment de l'usage, guidage vers Réglages en cas de refus
  \item La fenêtre de permission iOS ne passe qu'une fois : choisis « l'instant où on s'en sert », et écris clairement l'usage
  \item Trois disciplines de notification : ne diffuser que l'utile, pouvoir couper, ne pas harceler ; framework UserNotifications
  \item MVVM avancé : abstraction par protocole (testable), responsabilité unique, pas de surconception
  \item Dividendes de l'écosystème : widgets WidgetKit + raccourcis App Intents
  \item Quatre scénarios d'assistance par l'IA : coder, chercher la doc, écrire les tests, corriger les bugs ; méthode en septième partie (chapitres 37-47)
\end{itemize}

\end{tcbsummary}
